% -----------------------------------------------------------
\section{Firstfruits}
% -----------------------------------------------------------
	\label{FIRSTFRUITS}
	
% ----------------------
\subsection{See also}
% ----------------------

	\hyperref[END]{End}, \hyperref[INTERCESSION]{Intercession}, \hyperref[NEWBIRTH]{New Birth}, \hyperref[RECONCILIATION]{Reconciliation}, \hyperref[TEMPLE]{Temple},

% ----------------------
\subsection{1828 American Dictionary of the English Language} % *1828
% ----------------------
	\label{FIRSTFRUITS-1828}

	\begin{compactenum}
		\item[1] The fruit or produce first matured and collected in any season. Of these the Jews made an oblation to God, as an acknowledgment of his sovereign dominion.
		\item[2] The first profits of any thing. In the church of England, the profits of every spiritual benefice for the first year.
		\item[3] The first or earliest effect of any thing, in a good or bad sense; as the first-fruits of grace in the heart, or the first-fruits of vice.
	\end{compactenum}

% ----------------------
\subsection{Oxford English Dictionary} % *OED
% ----------------------
	\label{FIRSTFRUITS-OED}

	\begin{compactitem}
		\item[1] The first agricultural products of a season, esp. with reference to the custom of making offerings of these to God or the gods.
		\item[2] \textit{figurative} and in extended use.
		\item[2.a] In \textit{plural}. The earliest products or results of anything; the first products of a person's work or endeavour. Also (with singular reference): something representative of a greater quantity that is expected later.
		\item[2.b] In \textit{singular}. A first product or result; something representative of a greater quantity that is expected later.
	\end{compactitem}

% ----------------------
\subsection{Scriptural Usage} % *SCRIPTURES
% ----------------------
	\label{FIRSTFRUITS-SCRIPTURES}

	% -----
	\subsubsection{Old Testament} % *OT
	% -----
		\label{FIRSTFRUITS-OT}
	
		% --
		\paragraph{KJV}
		% --
			\label{FIRSTFRUITS-OT-KJV}
	
			\begin{tabular}{ll}
				Total occurrences in the OT & 26
			\end{tabular}
			
			\begin{compactdesc}
				\item[Exodus 23:16] And the feast of harvest, the firstfruits of thy labours, which thou hast sown in the field: and the feast of ingathering, which is in the end of the year, when thou hast gathered in thy labours out of the field.
				\item[Exodus 23:19] The first of the firstfruits of thy land thou shalt bring into the house of the LORD thy God. Thou shalt not seethe a kid in his mother's milk.
				\item[Leviticus 2:12] As for the oblation of the firstfruits, ye shall offer them unto the LORD: but they shall not be burnt on the altar for a sweet savour.
				\item[Numbers 18:12--15] All the best of the oil, and all the best of the wine, and of the wheat, the firstfruits of them which they shall offer unto the LORD, them have I given thee. 13 And whatsoever is first ripe in the land, which they shall bring unto the LORD, shall be thine; every one that is clean in thine house shall eat of it. 14 Every thing devoted in Israel shall be thine. 15 Every thing that openeth the matrix in all flesh, which they bring unto the LORD, whether it be of men or beasts, shall be thine: nevertheless the firstborn of man shalt thou surely redeem, and the firstling of unclean beasts shalt thou redeem.
				\item[2 Chronicles 31:5] And as soon as the commandment came abroad, the children of Israel brought in abundance the firstfruits of corn, wine, and oil, and honey, and of all the increase of the fields; and the tithe of all things brought they in abundantly.
				\item[Nehemiah 10:35--37] And to bring the firstfruits of our ground, and the firstfruits of all fruit of all trees, year by year, unto the house of the LORD 36 Also the firstborn of our sons, and of our cattle, as it is written in the law, and the firstlings of our herds and of our flocks, to bring to the house of our God, unto the priests that minister in the house of our God: 37 And that we should bring the firstfruits of our dough, and our offerings, and the fruit of all manner of trees, of wine and of oil, unto the priests, to the chambers of the house of our God; and the tithes of our ground unto the Levites, that the same Levites might have the tithes in all the cities of our tillage.
				\item[Proverbs 3:9] Honour the LORD with thy substance, and with the firstfruits of all thine increase:
				\item[Jeremiah 2:3] Israel was holiness unto the LORD, and the firstfruits of his increase: all that devour him shall offend; evil shall come upon them, saith the LORD.
				\item[Ezekiel 20:40] For in mine holy mountain, in the mountain of the height of Israel, saith the Lord GOD, there shall all the house of Israel, all of them in the land, serve me: there will I accept them, and there will I require your offerings, and the firstfruits of your oblations, with all your holy things.
				\item[Ezekiel 48:14] And they shall not sell of it, neither exchange, nor alienate the firstfruits of the land: for it is holy unto the LORD.
			\end{compactdesc}
			
		% --
		\paragraph{Hebrew Bible}
		% --
			\label{FIRSTFRUITS-HEBREW}
		
			%
			\subparagraph{\texorpdfstring{\he{b*ik*Uriym}}{}}
			%
				\label{BIKKURIM} % whatever is in step bible without periods
				
				Very important to note the connection of this word to \he{b*Akar}, ``to be born first, to bear early.'' Christ is spoken of as the ``firstfruits unto God'' in the BOM, and we ourselves are spoken of as the ``first-fruits of Christ unto God.'' I think these uses of ``firstfruits'' are related to \he{b*ik*Uriym}, which is related to \he{b*Akar}. See \hyperref[FIRSTFRUITS-UNTOGOD]{this study}.
			
				\begin{compactdesc}
					\item[HALOT 1204, PDF 583] \textbf{}
						\begin{compactitem}
							\item early fruits, first-fruits (of grapes, seed, fields, trees, wheat, and more generally as well; Exodus 23:16
						\end{compactitem}
					\item[BDB 1061, PDF 334] \textbf{}
						\begin{compactitem}
							\item first-fruits; the first of grain and fruit that ripened and was gathered and offered to God according to the ritual
						\end{compactitem}
					\item[DCH PDF 647] \textbf{}
						\begin{compactitem}
							\item firstfruits; first portion of gathered cereal and fruit crops, offered to God
						\end{compactitem}
				\end{compactdesc}
				
				\begin{compactdesc}
					\item[Some OT uses] \textbf{}
						\begin{compactdesc}
							\item[Exodus 23:16] firstfruits
							\item[Exodus 23:19] firstfruits
							\item[Numbers 18:13] whatsoever is first ripe
							\item[Nehemiah 10:36] firstfruits
						\end{compactdesc}
				\end{compactdesc}
				
			%
			\subparagraph{\texorpdfstring{\he{re'+siyt}}{}}
			%
				\label{RESHIT} % whatever is in step bible without periods
			
				\begin{compactdesc}
					\item[HALOT 8618, PDF 2465] \textbf{}
						\begin{compactitem}
							\item[1] what comes first, beginning
							\item[2] beginning, starting point
							\item[3] the first and best
							\item[3.b] the first and best part (position)
							\item[3.b.i.] [the first and best part] of the produce
							\item[4] in the context of ritual, first-fruit, choicest portion (Leviticus 2:12)
						\end{compactitem}
					\item[BDB 7225, PDF 2216] \textbf{}
						\begin{compactitem}
							\item beginning, chief
							\item[1.a] beginning
							\item[1.b] first of fruits (Leviticus 2:12) 
							\item[2] first, chief
						\end{compactitem}
					\item[DCH PDF 4538] \textbf{}
						\begin{compactitem}
							\item beginning, first
							\item[1] beginning
							\item[2] first (one), first (thing), firstfruits (Leviticus 2:12)
							\item[3] first, i.e. chief (one, thing, part), choice(st thing, part), choice produce
						\end{compactitem}
				\end{compactdesc}
				
				\begin{compactdesc}
					\item[Some OT uses] \textbf{}
						\begin{compactdesc}
							\item[Leviticus 2:12] firstfruits
							\item[Numbers 18:12] firstfruits
							\item[2 Chronicles 31:5] firstfruits
							\item[Proverbs 3:9] firstfruits
							\item[Jeremiah 2:3] firstfruits
							\item[Ezekiel 20:40] firstfruits
							\item[Ezekiel 48:14] firstfruits
						\end{compactdesc}
				\end{compactdesc}
			
		% % --
		% \paragraph{Septuagint}
		% % --
			% \label{FIRSTFRUITS-LXX}
		
			% %
			% \subparagraph{\texorpdfstring{\gr{sample word}}{}}
			% %
		
	% -----
	\subsubsection{New Testament} % *NT
	% -----
		\label{FIRSTFRUITS-NT}
	
		% --
		\paragraph{KJV}
		% --
			\label{FIRSTFRUITS-NT-KJV}
	
			\begin{tabular}{ll}
				Total occurrences in the NT & 8
			\end{tabular}
			
			\begin{compactdesc}
				\item[Romans 8:23] And not only they, but ourselves also, which have the firstfruits of the Spirit, even we ourselves groan within ourselves, waiting for the adoption, to wit, the redemption of our body.
				\item[Romans 11:16] For if the firstfruit be holy, the lump is also holy: and if the root be holy, so are the branches.
				\item[Romans 16:5] Likewise greet the church that is in their house. Salute my wellbeloved Epænetus, who is the firstfruits of Achaia unto Christ.
				\item[1 Corinthians 15:20] But now is Christ risen from the dead, and become the firstfruits of them that slept.
				\item[1 Corinthians 15:23] But every man in his own order: Christ the firstfruits; afterward they that are Christ's at his coming.
				\item[1 Corinthians 16:15] I beseech you, brethren, (ye know the house of Stephanas, that it is the firstfruits of Achaia, and that they have addicted themselves to the ministry of the saints,)
				\item[James 1:18] Of his own will begat he us with the word of truth, that we should be a kind of firstfruits of his creatures.
				\item[Revelation 14:4] These are they which were not defiled with women; for they are virgins. These are they which follow the Lamb whithersoever he goeth. These were redeemed from among men, being the firstfruits unto God and to the Lamb.
			\end{compactdesc}
			
		% --
		\paragraph{Greek New Testament}
		% --
			\label{FIRSTFRUITS-GREEK}
		
			%
			\subparagraph{\texorpdfstring{\gr{ἀπαρχή}}{}}
			%
				\label{APARCHE}
			
				\begin{compactdesc}
					\item[BDAG 85b] \textbf{}
						\begin{compactitem}
							\item[1] \textit{first fruits, first portion} of any kind
							\item[2] \textit{birth-certificate}
						\end{compactitem}
					\item[LSJ 180b, PDF 253] \textbf{}
						\begin{compactitem}
							\item[1] \textit{beginning of a sacrifice, primal offering}
							\item[2] \textit{firstlings} for sacrifice or offering, \textit{first-fruits}
							\item[7] \textit{birth-certificate} of a free person
						\end{compactitem}
				\end{compactdesc}
				
				\begin{compactdesc}
					\item[Some NT uses] \textbf{}
						\begin{compactdesc}
							\item[Romans 8:23] \gr{οὐ μόνον δέ, ἀλλὰ καὶ αὐτοὶ τὴν ἀπαρχὴν τοῦ πνεύματος ἔχοντες ἡμεῖς καὶ αὐτοὶ ἐν ἑαυτοῖς στενάζομεν, υἱοθεσίαν ἀπεκδεχόμενοι τὴν ἀπολύτρωσιν τοῦ σώματος ἡμῶν.}
							\item[Romans 11:16] \gr{εἰ δὲ ἡ ἀπαρχὴ ἁγία, καὶ τὸ φύραμα· καὶ εἰ ἡ ῥίζα ἁγία, καὶ οἱ κλάδοι.}
							\item[Romans 16:5] \gr{καὶ τὴν κατ᾽ οἶκον αὐτῶν ἐκκλησίαν. ἀσπάσασθε Ἐπαίνετον τὸν ἀγαπητόν μου, ὅς ἐστιν ἀπαρχὴ τῆς Ἀσίας εἰς Χριστόν.}
							\item[1 Corinthians 15:20] \gr{Νυνὶ δὲ Χριστὸς ἐγήγερται ἐκ νεκρῶν, ἀπαρχὴ τῶν κεκοιμημένων.}
							\item[1 Corinthians 15:23] \gr{ἕκαστος δὲ ἐν τῷ ἰδίῳ τάγματι· ἀπαρχὴ Χριστός, ἔπειτα οἱ τοῦ Χριστοῦ ἐν τῇ παρουσίᾳ αὐτοῦ·}
							\item[1 Corinthians 16:15] \gr{Παρακαλῶ δὲ ὑμᾶς, ἀδελφοί· οἴδατε τὴν οἰκίαν Στεφανᾶ, ὅτι ἐστὶν ἀπαρχὴ τῆς Ἀχαΐας καὶ εἰς διακονίαν τοῖς ἁγίοις ἔταξαν ἑαυτούς·}
							\item[2 Thessalonians 2:13] \gr{Ἡμεῖς δὲ ὀφείλομεν εὐχαριστεῖν τῷ θεῷ πάντοτε περὶ ὑμῶν, ἀδελφοὶ ἠγαπημένοι ὑπὸ κυρίου, ὅτι εἵλατο ὑμᾶς ὁ θεὸς ἀπαρχὴν εἰς σωτηρίαν ἐν ἁγιασμῷ πνεύματος καὶ πίστει ἀληθείας,} \textbf{from the beginning} (The KJV translation may be wrong. It takes the word as \gr{ἀπ᾽ ἀρχῆς}, which it apparently is in some manuscripts. But the NRSV translation of this verse, for example, reads ``But we must always give thanks to God for you, brothers and sisters beloved by the Lord, because God chose you as the first fruits for salvation through sanctification by the Spirit and through belief in the truth.'' See \hyperref[2THESSALONIANS-2-13-NOTE-beginning]{note at 2 Thessalonians 2:13}.)
							\item[James 1:18] \gr{βουληθεὶς ἀπεκύησεν ἡμᾶς λόγῳ ἀληθείας, εἰς τὸ εἶναι ἡμᾶς ἀπαρχήν τινα τῶν αὐτοῦ κτισμάτων.}
							\item[Revelation 14:4] \gr{οὗτοί εἰσιν οἳ μετὰ γυναικῶν οὐκ ἐμολύνθησαν, παρθένοι γάρ εἰσιν· οὗτοι οἱ ἀκολουθοῦντες τῷ ἀρνίῳ ὅπου ἂν ὑπάγῃ· οὗτοι ἠγοράσθησαν ἀπὸ τῶν ἀνθρώπων ἀπαρχὴ τῷ θεῷ καὶ τῷ ἀρνίῳ,}
						\end{compactdesc}
				\end{compactdesc}
		
	% -----
	\subsubsection{Book of Mormon} % *BOM
	% -----
		\label{FIRSTFRUITS-BOM}
	
		\begin{tabular}{ll}
			Total occurrences in the BOM & 2
		\end{tabular}
		
		\begin{compactdesc}
			\item[2 Nephi 2:9] Wherefore, he is the firstfruits unto God, inasmuch as he shall make intercession for all the children of men; and they that believe in him shall be saved. (See notes at \hyperref[2NEPHI-2-9]{2 Nephi 2:9}.)
			\item[Jacob 4:11] Wherefore, beloved brethren, be reconciled unto him through the atonement of Christ, his Only Begotten Son, and ye may obtain a resurrection, according to the power of the resurrection which is in Christ, and be presented as the first-fruits of Christ unto God, having faith, and obtained a good hope of glory in him before he manifesteth himself in the flesh.
		\end{compactdesc}
		
	% -----
	\subsubsection{Doctrine and Covenants} % *DC
	% -----
		\label{FIRSTFRUITS-DC}
	
		\begin{tabular}{ll}
			Total occurrences in the DC & 0
		\end{tabular}
		
		% \begin{compactdesc}
			% \item[]
		% \end{compactdesc}
		
	% -----
	\subsubsection{Pearl of Great Price} % *PGP
	% -----
		\label{FIRSTFRUITS-PGP}
	
		\begin{tabular}{ll}
			Total occurrences in the PGP & 0
		\end{tabular}
		
		% \begin{compactdesc}
			% \item[]
		% \end{compactdesc}
		
% ----------------------
\subsection{Key Phrases} % *KEYPHRASES *PHRASES
% ----------------------
	\label{FIRSTFRUITS-PHRASES}

	\begin{compactdesc}
		\item[firstfruits unto God] Revelation 14:4; 2 Nephi 2:9; Jacob 4:11
	\end{compactdesc}
	
% % ----------------------
% \subsection{Bible Dictionaries} % *DICTIONARIES
% % ----------------------
	% \label{FIRSTFRUITS-BD}

% % ----------------------
% \subsection{Restoration Usage} % *RESTORATION
% % ----------------------
	% \label{FIRSTFRUITS-RESTORATION}

	% % -----
	% \subsubsection{Joseph Smith} % JS
	% % -----
		% \label{FIRSTFRUITS-JS}
	
		% \begin{compactdesc}
			% \item[DATE/SOURCE]
		% \end{compactdesc}
	
	% % -----
	% \subsubsection{Brigham Young} % BY
	% % -----
		% \label{FIRSTFRUITS-BY}
	
		% \begin{compactdesc}
			% \item[DATE/SOURCE]
		% \end{compactdesc}
	
% ----------------------
\subsection{Commentary and Studies} % *COMMENTARY *STUDIES
% ----------------------
	\label{FIRSTFRUITS-COMMENTARY}
	
	% % -----
	% \subsubsection{Key Points}
	% % -----
		% \label{FIRSTFRUITS-KEYS}
	
		% \begin{compactitem}
			% \item
		% \end{compactitem}
		
	% -----
	\subsubsection{Firstfruits unto God}
	% -----
		\label{FIRSTFRUITS-UNTOGOD}
	
		Both Christ and we, as God's children, are spoken of as ``firstfruits'' in various places in the scriptures:
		
		\begin{description}
			\item[Christ] 1 Corinthians 15:20, 23; 2 Nephi 2:9
			\item[God's children] Romans 16:5; James 1:18; Revelation 14:4; Jacob 4:11
		\end{description}
		
		If you think of ``firstfruits'' as the first mature product of some process, then I think it is possible to see two different processes in play here, both of which produce ``firstfruits'':
		
		\begin{enumerate}
			\item The process Christ is currently undergoing, or underwent and finished with the atonement and resurrection;
			\item The process we are currently undergoing, this stage of which will end with our exaltation in the celestial kingdom.
		\end{enumerate}
		
		I like the idea of separating the processes because, even though we and Christ are both referred to as the ``firstfruits unto God,'' it is clear that we have different roles to play, and very different things are required of us. I think this is an indication of the two different processes: Christ is the firstfruits of his process ``inasmuch as he shall make intercession for all the children of men'' (2 Nephi 2:9), and we are the firstfruits of our process through the redemption, reconciliation, and new birth (Revelation 14:4, Jacob 4:11, James 1:18).
		
		I think these two different processes also speak to the different sources of power: Christ's power is directly through Heavenly Father (\hyperref[JOHN-5-26]{John 5:26}, though that isn't the best verse for it, I think), and ours is through Christ (\hyperref[DC35-2]{DC 35:2}).